\documentclass[11pt,a4paper]{article}
\usepackage[utf8]{inputenc}
\usepackage[T1]{fontenc}

\usepackage{amsmath,amssymb}
\usepackage[margin=2.3cm]{geometry}
\usepackage{enumitem}
\usepackage{xcolor}
\usepackage{tcolorbox}
\usepackage{booktabs}

\definecolor{navy}{HTML}{2B2A6B}
\definecolor{coral}{HTML}{C2406F}
\setlist{itemsep=2pt,topsep=4pt}
\setlength{\parindent}{0pt}\setlength{\parskip}{6pt}
\newtcolorbox{aviso}{colback=yellow!10,colframe=navy,boxrule=.6pt,arc=3pt,left=6pt,right=6pt,top=4pt,bottom=4pt}
\newcommand{\seccion}[1]{\vspace{4pt}{\large\bfseries\color{navy}#1}\par\vspace{-2pt}}
\begin{document}
{\LARGE\bfseries\color{navy}Guía de estilo y plantilla común de estación}\\[2pt]
{\color{coral}\bfseries PACE · UACh · Curso de nivelación en matemática}\\[2pt]
{\small Borrador para revisión del equipo docente · noviembre de 2026}

\begin{aviso}
Esta guía \textbf{orienta}, no obliga. Cada docente decide cómo explicar sus contenidos y cuánto dura su unidad. Lo que se comparte son la organización visual, los términos de la interfaz y unos criterios mínimos de calidad. Los puntos marcados con \textbf{(a acordar)} son propuestas que el equipo debe confirmar.
\end{aviso}

\seccion{1. Estructura de una estación}
Cada estación tiene tres pasos, siempre en este orden: \textbf{Aprende}, \textbf{Practica} y \textbf{Aplica}. En Aprende se sugiere esta secuencia, que en el editor corresponde al bloque «Aprende visual»:
\begin{enumerate}
 \item \textbf{Qué aprenderás.} Un objetivo observable, en una o dos oraciones («Al terminar podrás resolver ecuaciones lineales con paréntesis»).
 \item \textbf{Comprende la idea.} La explicación y su fundamento: por qué funciona, no solo cómo se hace.
 \item \textbf{Ejemplo guiado.} Un ejemplo resuelto en pasos cortos. Cada paso dice qué se hace y por qué, y muestra la expresión matemática.
 \item \textbf{Presta atención.} Un error frecuente o una condición que conviene vigilar, con un ejemplo breve.
 \item \textbf{Para profundizar (opcional).} Ejemplos adicionales que el estudiante abre si los quiere.
\end{enumerate}
Ningún elemento debe ser un párrafo largo. Si una explicación necesita más de unas pocas oraciones, conviene dividirla en pasos.

\seccion{2. Lenguaje}
\begin{itemize}
 \item Claro, directo y preciso. Se trata al estudiante de «tú», como a una persona adulta, sin tono infantil.
 \item Evitar reglas sin fundamento. Por ejemplo, no escribir «el $3$ pasa sumando»; escribir «restamos $3$ en ambos lados».
 \item Una idea por oración. Definir cada término la primera vez que aparece.
 \item Los enunciados de Aplica usan contextos de ciencias (física, química, biología, ciencias de la Tierra) con datos verosímiles.
\end{itemize}

\seccion{3. Notación matemática \textnormal{\small (a acordar)}}
Propuestas iniciales para que las tres unidades se lean igual:
\begin{center}
\begin{tabular}{@{}ll@{}}
\toprule
\textbf{Tema} & \textbf{Propuesta}\\
\midrule
Decimales & coma decimal: $0{,}25$\\
Multiplicación & $2x$ o $2\cdot 3$; no usar «$\times$» ni «$*$»\\
Fracciones & $\dfrac{a}{b}$ en pantalla; en línea, $\tfrac{a}{b}$ solo si es sencilla\\
Potencias y raíces & $x^{2}$, $\sqrt{x}$, $\sqrt[3]{x}$\\
Intervalos y conjuntos & $[a,b]$, $(a,b)$, $\{x\in\mathbb{R} : x>2\}$\\
Aproximaciones & indicar siempre cuántas cifras se conservan: $\pi\approx 3{,}14$\\
Fórmula resolvente & $x=\dfrac{-b\pm\sqrt{b^{2}-4ac}}{2a}$\\
\bottomrule
\end{tabular}
\end{center}
En el editor, las fórmulas van entre \verb|$...$| (en el texto) o \verb|$$...$$| (centradas). Para moneda se escribe «CLP 3000», sin usar \verb|$|.

\seccion{4. Términos de la interfaz}
\begin{itemize}
 \item Los tres pasos se llaman \textbf{Aprende}, \textbf{Practica} y \textbf{Aplica}, y la evaluación de cierre se llama \textbf{META FINAL}.
 \item Un bloque agrupa estaciones; una unidad agrupa bloques. El curso no se aprueba ni se reprueba: se \textbf{completa}.
 \item Cuando el banco no garantiza preguntas nuevas se dice «\textbf{Practicar otra ronda}», no «nuevas preguntas».
 \item Se distingue siempre el \textbf{avance acumulado} (lo que queda guardado) del resultado de la ronda actual.
 \item La mascota es el chucao, que cambia de etapa según el avance de la unidad.
\end{itemize}

\seccion{5. Retroalimentación}
\begin{itemize}
 \item Si la respuesta es correcta: confirmar y explicar brevemente por qué.
 \item Si es incorrecta: no revelar la respuesta; indicar qué revisar y animar a intentarlo de nuevo.
 \item Las alternativas incorrectas se diseñan a partir de errores reales y, cuando se pueda, cada una tiene su propia pista.
\end{itemize}

\seccion{6. Revisión mínima de una estación}
\begin{enumerate}[label=$\square$]
 \item El objetivo es claro y observable.
 \item Los ejemplos son matemáticamente correctos (otra persona los resolvió).
 \item La notación es consistente con la sección 3.
 \item La retroalimentación se entiende sin ayuda.
 \item Se ve bien en un teléfono, sin desplazamiento horizontal.
 \item La exportación funciona y el SCORM carga en SIVEDUC.
 \item Practica + Aplica + Aprende suman como máximo 24 actividades por estudiante (el editor avisa).
\end{enumerate}

\seccion{7. Cómo se usa en el editor}
Abre el paquete de tu unidad, elige una estación y, en «Editar estación», activa «Aprende visual». Completa objetivo, idea, pasos y «presta atención»; la vista previa de la derecha muestra el resultado. La estación de ejemplo es la 3 de la Unidad III, «Paréntesis y fracciones», que ya tiene todos los elementos.
\end{document}
